\documentclass[11pt,a4paper]{article}
\usepackage[T1]{fontenc}
\usepackage[utf8]{inputenc}
\usepackage[margin=2.2cm]{geometry}
\usepackage{booktabs}
\usepackage{array}
\usepackage{xcolor}
\usepackage{fancyvrb}
\usepackage{parskip}
\usepackage[hidelinks]{hyperref}
\usepackage{titlesec}

\definecolor{good}{HTML}{157A3C}
\definecolor{bad}{HTML}{B3261E}
\definecolor{boxbg}{HTML}{F2F2F5}
\newcommand{\good}[1]{\textcolor{good}{\textbf{#1}}}
\newcommand{\bad}[1]{\textcolor{bad}{\textbf{#1}}}

\titleformat{\section}{\large\bfseries}{\thesection}{0.6em}{}
\titleformat{\subsection}{\normalsize\bfseries}{\thesubsection}{0.6em}{}
\setlength{\tabcolsep}{5pt}
\renewcommand{\arraystretch}{1.15}

\title{\textbf{TRELLIS.2 Printability Fine-Tuning}\\[2pt]\large Weekly Report}
\author{Branch \texttt{flow-dpo-training}\\Apple M4, 32\,GB unified memory}
\date{2026-07-31 to 2026-08-06 \quad$\cdot$\quad 17 commits}

\begin{document}
\maketitle

\begin{abstract}
Many post-processing libraries and methods were tried to improve quality and printability, adding
different options for post-processing. \textbf{manifold3d} was found to be the best for 3D printing, and the
original trimesh pipeline was kept as well. Further research was done on the best way to decrease
non-manifold edges of the models, and the TriFlow paper on generating artist-like 3D models was read.

Integrating DreamDPO concepts was also tried, with no success. The DreamDPO work for diffusion
models was read, but the preference discriminator was simulated as a score rather than a VLM, and the
results showed no gains --- it always chose the reference path for the flow model. That, together with
a lack of expertise in this area of the theory, led to the decision to work instead on fine-tuning the
SLat flow model using LoRA methods.

This was tried on a small sample size of 300 curated 3D models from Thingi10K. Initial training
showed gains, with significantly less non-manifold edges for raw results, without loss in detail. More
latent data is being collected for training.

\end{abstract}
\section{Executive summary}
We retired the inference-time DPO steering approach, established that supervised fine-tuning is the
better instrument for this problem, built the full data and training pipeline, and completed a
1,170-step LoRA fine-tune of TRELLIS.2's shape-SLat flow model on 300 clean Thingi10K models.

\textbf{The fine-tune works on the axes it can reach.} Measured on exported meshes, held-out set,
paired sampling noise:

\begin{center}\small
\begin{tabular}{@{}lrrrr@{}}
\toprule
\textbf{Model} & \textbf{L\_OH (overhang)} & \textbf{L\_Th (thin walls)} & \textbf{L\_Topo} & \textbf{R\_Detail} \\
\midrule
Base model & 0.2291 & 0.1133 & 0.1871 & 0.4316 \\
step 750 & 0.2183 (-5\%) & \good{0.0152 (-87\%)} & 0.0500 (-73\%) & 0.4578 (+6\%) \\
step 900 & 0.2255 (-2\%) & 0.0299 (-74\%) & 0.0635 (-66\%) & 0.4411 (+2\%) \\
\textbf{step 1050} & 0.2179 (-5\%) & \good{0.0159 (-86\%)} & \good{0.0340 (-82\%)} & 0.4358 (+1\%) \\
\bottomrule
\end{tabular}
\end{center}
Step 1050 beats the base model on thin-wall penalty for \textbf{6 of 6} held-out objects. Image
fidelity did \textit{not} degrade (0.897 $\rightarrow$ 0.908--0.913), and surface detail slightly increased,
so this is not printability bought by smoothing the model into generic blobs.

\begin{center}\colorbox{boxbg}{\parbox{0.95\linewidth}{\vspace{3pt}\textbf{Three caveats stated up front.}
(1) Overhang never moved on any checkpoint.
(2) Watertightness is unchanged and cannot be changed by this training --- it is a decoder property.
(3) The run diverged at \textasciitilde{}step 1100 due to a missing LR schedule; the last 70 steps are worthless and
the final model is \textit{not} the best model.\vspace{3pt}}}\end{center}
\section{What was discovered}
\subsection{The retired approach failed for a diagnosable reason}
Rather than using a VLM, the preference discriminator was simulated as a score. It produced no
gains, and always chose the reference path for the flow model. The measurements below say the same
thing from the other side.

Inference-time DPO steering produced a gradient contribution of \textbf{-0.0027 $\pm$ 0.0078} over
17 scored forks --- indistinguishable from zero. Root cause: the only differentiable scalar available
(\texttt{slat\_detail\_proxy}) agreed with the judge's within-fork argmax \textbf{5.9\%} of the time against
a 33\% chance line, with mean Spearman -0.265. The optimiser worked; its objective was
near-orthogonal to the goal.

\subsection{The "differentiability wall" did not exist}
The premise that mesh extraction blocks gradients is false. TRELLIS.2's decoder is dual-grid and
regresses vertex positions smoothly (\texttt{vertices = (1+2$\cdot$margin)$\cdot$sigmoid(...) - margin});
only connectivity is boolean. Six weeks of design rested on an unverified assertion that took one
inspection of two source files to refute.

\subsection{The VAE bounds what any flow-model training can achieve}
Encoding a perfectly watertight mesh and decoding it straight back:

\begin{center}\small
\begin{tabular}{@{}ll@{}}
\toprule
\textbf{Property} & \textbf{Result} \\
\midrule
Shape fidelity & \good{Preserved --- Chamfer 0.0052, p95 0.0101, bounding box exact} \\
Watertightness & \bad{0 of 2 preserved} \\
Defect rate (full mesh) & 0.086\% open edges, 0.095\% non-manifold \\
Components & 953, but largest is 99.85\% of faces --- and still not watertight alone \\
\bottomrule
\end{tabular}
\end{center}
\textbf{Consequence:} geometry-carried printability (thickness, overhang, bulk shape) is reachable by
training the flow model; watertightness is not, and belongs to post-processing. This single 20-minute
measurement redirected the entire project.

\subsection{The judge inflates topology defects \textasciitilde{}50x}
\texttt{geometric\_judge} decimates before scoring. Measured on the same mesh: \textbf{0.086\%} open-edge
rate on the full mesh versus \textbf{4.25\%} after the judge's preprocessing. Every historical
\texttt{L\_Topo} figure is therefore dominated by decimation artifacts rather than real decoder
defects. Any topology claim must be measured on the undecimated mesh.

\subsection{Measurement noise nearly made the whole run uninterpretable}
Two evaluations of the \textit{same} base model gave L\_OH 0.2889 and 0.3871 --- a 34\% swing from
sampling noise alone, larger than any plausible training effect. Fixed by seeding sampling noise per
held-out object and reusing it at every evaluation (common random numbers). Consecutive evals are now
bit-identical, so any movement is attributable to the model.

\subsection{Silent failures dominated the engineering cost}
Every significant delay this week was a silent failure, not a crash:

\begin{center}\small
\begin{tabular}{@{}lll@{}}
\toprule
\textbf{Symptom} & \textbf{Actual cause} & \textbf{Cost} \\
\midrule
Training step took 464s & No activation checkpointing $\rightarrow$ 38~GB retained on a 32~GB machine $\rightarrow$ swap thrashing (24.4~GB of 25.6~GB swap in use). Did not OOM, just degraded. & 18x slowdown \\
Baseline eval returned \texttt{n\_eval=0} & Three stacked bugs, masked by \texttt{except: continue} & 2 invalidated launches \\
Eval decoder on CPU & \texttt{low\_vram} is device placement, not a memory hint; clearing it stranded the decoder & 1 launch \\
Sample-evolution panel empty & A patch silently failed to match; the call site was never added & whole run, cosmetic \\
Loss chart nonsense & Dataset build logged token counts into the \texttt{loss} field & misleading chart \\
\bottomrule
\end{tabular}
\end{center}
{\small The single highest-leverage change all week was replacing \texttt{except Exception: continue}
with recorded, printed failures. It converted two blind retries into one-step diagnoses.}

\subsection{Post-processing was addressed before generation was}
Many post-processing libraries and methods were tried to improve quality and printability, with
different options added to the pipeline. \textbf{manifold3d} was found to be the best for 3D printing and
was added, with the original trimesh pipeline kept alongside it. Separate research was done on the
best way to decrease non-manifold edges, and the TriFlow paper on generating artist-like 3D models
was read.

This matters for reading the rest of this report: the fine-tuning results quoted here are on
\textbf{raw} decoder output, with no post-processing applied. The reduction in non-manifold edges is
therefore a gain made \textit{before} manifold3d or trimesh repair runs, not instead of them. Both remain
in the pipeline, and the round-trip result in \S{}2.3 is why --- watertightness is a decoder property
that fine-tuning cannot reach.

\section{What was built}
\begin{center}\small
\begin{tabular}{@{}lll@{}}
\toprule
\textbf{Component} & \textbf{File} & \textbf{Note} \\
\midrule
CPU mesh $\rightarrow$ SLat encoder & \texttt{backends/dual\_grid\_cpu.py} & o-voxel ships CUDA-only; rebuilt the pure-C++ voxelizer for Apple Silicon \\
Offscreen renderer & \texttt{trellis\_core/render\_mesh.py} & moderngl, Metal-backed, headless; camera fitted per-mesh to bounding sphere \\
Dataset pipeline & \texttt{trellis\_core/dataset\_build.py} & 300 models, 3.2 GPU-hours, 0 failures \\
LoRA & \texttt{trellis\_core/lora.py} & 210 modules, 18.68M params (1.4\% of 1.29B) \\
SFT trainer & \texttt{trellis\_core/sft\_train.py} & Hyperparameters from TRELLIS.2's own config \\
Live dashboard & \texttt{devlab/train\_server.py} & Crash-independent; phone access over Tailscale \\
VAE round-trip test & \texttt{trellis\_core/vae\_roundtrip.py} & Established the ceiling \\
Mesh exporter & \texttt{trellis\_core/export\_samples.py} & Reproduces eval exactly (same seeds, same CFG) \\
Topology test & \texttt{trellis\_core/topology\_test.py} & Paired, undecimated, multi-seed \\
\bottomrule
\end{tabular}
\end{center}
\subsection{Dataset}
Thingi10K via its HuggingFace mirror. The corpus is deliberately a record of real-world
3D-printing messiness --- its own figures are 50\% non-solid, 45\% self-intersecting, 26\%
multi-component, 22\% non-manifold. Filtering to \textit{Closed + edge/vertex manifold + single component +
PWN + no degenerate faces + zero self-intersections} leaves \textbf{3,935 of 10,000}. We used 300
(270 train / 30 held-out); median latent 2,233 voxels; 666~MB.

\subsection{Training run}
\begin{Verbatim}[fontsize=\small,frame=leftline,framerule=1pt,xleftmargin=8pt]
run id        sft-1200-20260806-0119
steps         1,170 of 1,200 (stopped after divergence)
epochs        8.7 over 270 examples
GPU time      9.96 hours, mean 30.7 s/step
LoRA          rank 16, alpha 32, lr 1e-4 constant, AdamW wd 0.01
schedule      uniform t, p_uncond 0.1 (both from TRELLIS.2's official config)
evals         10, every 150 steps, 6 held-out objects, common random numbers
\end{Verbatim}
\section{The divergence}
\begin{center}\small
\begin{tabular}{@{}lrrr@{}}
\toprule
\textbf{Steps} & \textbf{Mean loss} & \textbf{Max loss} & \textbf{Grad norm} \\
\midrule
1--1100 & \textasciitilde{}0.16 & 0.80 & 0.03--0.06 \\
1101--1170 & \bad{3.48} & \bad{57.29} & \bad{332.25} \\
\bottomrule
\end{tabular}
\end{center}
Held-out loss went 0.0859 $\rightarrow$ 10.29; image similarity 0.908 $\rightarrow$ 0.202. The model was destroyed.
Two causes, both mine: \textbf{no learning-rate decay} (constant 1e-4 for 8.7 epochs), and gradient
clipping that bounded each update but could not prevent AdamW's moments accumulating through a
sustained blow-up. Checkpoints 750/900/1050 were preserved before stopping, because a clean stop
writes a final checkpoint and \texttt{prune\_checkpoints(keep=3)} would have deleted step 750.

\section{How to log improvements numerically}
Recommended standard for every future comparison. The first three are non-negotiable given what we
measured this week.

\subsection{Always paired, never pooled}
Fix the sampling noise per held-out object and reuse it across every model being compared. Report
\textit{per-object deltas} and a win count, not two means. Justification: 96.8\% of judge-score variance
is between-object, and a pooled metric previously scored raw vertex count at r=0.959 while a perfect
within-fork ranker scored 0.180.

\subsection{Always against a baseline}
Every run begins with a step-0 evaluation of the base model with adapters disabled. A printability
number without the base value on the same axes cannot indicate whether training helped.

\subsection{Always report fidelity beside printability}
Minimum pair: a printability metric and \texttt{image\_similarity} (silhouette IoU against the
conditioning render). A printability gain with fidelity falling is a regression. Precedent: a 2-D
experiment tripled its reward while collapsing all 3,000 samples onto one geometry.

\subsection{The metric set}
\begin{center}\small
\begin{tabular}{@{}llll@{}}
\toprule
\textbf{Metric} & \textbf{Source} & \textbf{Direction} & \textbf{Notes} \\
\midrule
\texttt{L\_Th} & judge & lower & Responds strongly. Primary success axis. \\
\texttt{nonmanifold\_rate}, \texttt{open\_rate} & full mesh & lower & \textbf{Must} be computed undecimated (see 2.4) \\
\texttt{L\_OH} & judge & lower & Has not responded; also only \textasciitilde{}5\% of it is genuine support cost \\
\texttt{R\_Detail} & judge & context & Fidelity term, not printability --- no better/worse verdict \\
\texttt{image\_similarity} & silhouette IoU & higher & Prompt-adherence guard \\
\texttt{sample\_dispersion} & latent means & flat & Collapse guard \\
\texttt{watertight\_rate} & full mesh & control & Expected flat; a change indicates a bug, not progress \\
support volume & \textit{slicer, not yet built} & lower & The only metric outside the training reward \\
\bottomrule
\end{tabular}
\end{center}
\subsection{The acceptance bar}
Improvement over the base model is \textit{not} the bar. The bar is:

\begin{Verbatim}[fontsize=\small,frame=leftline,framerule=1pt,xleftmargin=8pt]
capture = (tuned_1_sample - base_1_sample) / (base_best_of_N - base_1_sample)
\end{Verbatim}
A single sample from the tuned model must recover a meaningful fraction of what N-fold rejection
sampling buys from the base model at matched compute. \textbf{This test has not yet been run.}

\section{Next steps}
\begin{enumerate}\itemsep3pt
  \item \textbf{Finish the topology consistency test} (running: 10 objects $\times$ 3 seeds $\times$ 4 models, paired,
undecimated). Answers whether step 1050's cleaner topology is consistent or a lucky draw.
  \item \textbf{Run the best-of-N comparison} (\S{}5.5). Until this exists we do not know whether the fine-tune
beats cheap rejection sampling --- the only comparison that justifies the work.
  \item \textbf{Fix the training schedule and re-run.} Cosine LR decay; a skip-step rule that \textit{drops} a batch
whose gradient norm exceeds a multiple of the running median rather than clipping it; keep every
checkpoint. The metrics were still improving when the run broke, so the ceiling was never observed.
  \item \textbf{Repair the judge} before it is trusted further: make the thickness raycast deterministic
(test-retest reliability was measured at 0.208, capping any correlation at |r| $\leq$ 0.456), and add a
support raycast to the overhang term (currently rank-correlates -0.432 with true unsupported area).
  \item \textbf{Scale the dataset.} 270 training objects is the binding constraint, not step count; 3,935 clean
models are available and prep runs at \textasciitilde{}40~s/model.
  \item \textbf{Add a slicer-based support-volume metric} as the held-out physical measure that was never part
of the training reward.
\end{enumerate}
\section{Deliverables on disk}
\begin{Verbatim}[fontsize=\small,frame=leftline,framerule=1pt,xleftmargin=8pt]
runs/sft-1200-20260806-0119/
    ckpt_keep/step_{750,900,1050}/adapter.pt   preserved LoRA adapters (643 MB)
    meshes/step_{0,750,900,1050}/              24 GLB + 24 STL, raw, no post-processing
    exports/                                   42 PNG frames + 6 progression strips
    metrics.jsonl                              1,170 train + 10 eval records
    topology_test.jsonl                        in progress
data/thingi10k_sft/                            300 latents + images + conditioning
\end{Verbatim}
{\small All meshes are raw decoder output --- no repair, no decimation, no solidify. Each object
uses identical conditioning and identical sampling noise across all four models, so differences are
attributable to the model alone.}

\end{document}
